\documentclass[11pt]{article}
\usepackage[a4paper,margin=27mm]{geometry}
\usepackage{amsmath,amssymb,amsthm,booktabs,hyperref}
\usepackage{xurl,microtype}
\setlength{\emergencystretch}{3em}
\newtheorem{theorem}{Theorem}
\newtheorem{lemma}[theorem]{Lemma}
\newtheorem{proposition}[theorem]{Proposition}
\newcommand{\Q}{\mathbb Q}\newcommand{\Z}{\mathbb Z}
\newcommand{\Tr}{\operatorname{Tr}}\newcommand{\sqf}{\operatorname{sqf}}
\title{Product poles and lower degrees in the few-variable constructions}
\author{\(\Z\Q\) Diophantine Project: audit and continuation}
\date{23 November 2026}
\begin{document}\maketitle
\begin{abstract}
We integrate the square-denominator criterion of Bhargava--Burungale,
arXiv:2611.16130v2, into the current finite-algebra constructions. Splitting
off an already available square factor before using their product equation
makes the radicand linear in the actual detector input. A class-seven
rescaling of the four-contact scalar makes the general sources compatible
with the published converse. Complete exact specifications give degree
bounds $48\,318\,382\,080$, $183\,368\,679\,424$ and $3\,261\,569$
in respectively 25, 27 and 11 total variables; the last construction
assumes ordinary rational abc. The 27-variable parameter remains $L=143$.
The other sufficient parameters remain ineffective. The two weak-abc
eleven-variable routes also improve, to degree bounds $16\,460\,929$
and $9\,808\,001$. All non-pole arithmetic
premises of the source constructions are expressly retained. These are
degree improvements, not new numerical variable-count records.
\end{abstract}

\paragraph{Attribution.}
The square-denominator theorem and its congruent-number-curve proof are
due to Bhargava and Burungale \cite{BB}. The source positive-factor,
terminal-root and projective encodings are from Gabe Conant's notes with
Astra~6.2, audited in v78. The point charts, contact-height reductions and
four-contact scalar are earlier project work. The normalized product
transfer, class-seven scalar, and present degree ledgers are subsequent
deductions. Dated sources and their hashes are preserved.

\section{Precise scope}
The base constructions are the two v78 papers \cite{scalar,twelve}.
Their retained inputs include the normalized elliptic group over
$\Q(\sqrt2)$, the appropriate trace or effective abscissa-sum height
theorem, the modular contact-height estimate or ordinary abc, and the
compatible simultaneous prime-pattern and index-refinement assertions.
The present proof replaces the old native-pole converse by
\cite[Lemma 2.2]{BB}; it does not establish those other inputs independently.
In particular it does not identify the paper's Green--Tao prime-progressions
lemma with the stronger simultaneous pattern used in our current sources.

Every coordinate count includes the input $X$. All polynomials below
are fully specified by finite multiplication algebras and determinants,
not by an unevaluated existence of coefficients. For the ineffective
families the fixed parameter is specified symbolically and the displayed
small member is only an algebraic fixture.

\begin{theorem}\label{main}
Under the corresponding retained premises, the following integer-polynomial
families have rational input projection exactly $\Z$.
\begin{center}\small
\begin{tabular}{lrrr}
\toprule
Route&Total variables&Degree at most&Support less than\\\midrule
General, sufficiently large odd $L$&25&$48\,318\,382\,080$&$2^{793}$\\
Fully specified, $L=143$&27&$183\,368\,679\,424$&$2^{906}$\\
Ordinary abc, sufficiently large odd $B$&11&$3\,261\,569$&$2^{202}$\\\bottomrule
\end{tabular}
\end{center}
The fully specified member has coefficient height less than $2^{2^{45}}$.
At the stored $L=143$ general fixture and $B=3$ abc fixture the coefficient
one-norm bounds are respectively
$2^{6235064631296}$ and $2^{12940032}$; neither fixture is asserted to be
a defining member. The degree and support bounds are uniform in the
admissible fixed parameters. Dividing by integer content makes each
polynomial primitive without increasing these bounds.
\end{theorem}

\section{Normalize before forming the product pole}
Call $e\in\Q^\times$ prepared if it has square reduced denominator and
\[
 0<e<1,\qquad v_2(e)\ge3,\qquad\sqf(e)\equiv7\pmod8.
\]
The final two conditions imply an even $v_2(e)\ge4$.

\begin{lemma}\label{pole}
Let $e=dk^2$ with $d,k\in\Q^\times$. Then the predicate
\begin{equation}\label{product}
 Y^2=(1-e)(v^4-d^2)(s^4-d^2),\qquad v,s,Y\in\Q,
\end{equation}
implies that $e$ has square reduced denominator. If $e$ is prepared,
the predicate has solutions, and its rational solutions project to a
Zariski-dense subset of the $(v,s)$-plane. Terminal zero roots are permitted.
\end{lemma}
\begin{proof}
Substitute $(s_0,t_0,z_0)=(kv,ks,k^4Y)$ in BB's equation
\[
 z_0^2=(1-e)(s_0^4-e^2)(t_0^4-e^2).
\]
Both quartic factors acquire the factor $k^4$, so this is an exact
equivalence. For soundness, if $v_p(e)<0$ is odd, then $v_p(1-e)$ is
odd and each quartic has finite even valuation, since the competing
valuations $4v_p(s_0)$ and $2v_p(e)$ cannot agree. Their product is
nonzero of odd valuation. This works at two as well as at odd primes.

For completeness put $j=\sqf(e)$ and $h=\sqf(1-e)$. BB Theorem~2.3
supplies points of prescribed abscissa squareclass on $E_n$ and $E_{nh}$;
when $h=1$ one may use the same positive-rank curve in both positions.
Odd multiples of either non-torsion point keep its Kummer abscissa class.
The nonconstant maps used in BB's proof of Lemma~2.2 therefore give
infinitely many rational first and second parameters, independently.
Every pair has the required product squareclass. The product of two
infinite subsets of $\Q$ is Zariski dense in $\mathbb A^2$; scaling by
$k^{-1}$ preserves this property. Consequently finitely many nonzero
polynomial guards can be avoided. No uniform statement is made for a
guard that becomes identically zero after a proposed parameter identification.
\end{proof}

The crucial algebraic feature is linear dependence on $e$ in
\eqref{product}. The preceding native equation had radicand
\[
 d(d+v^2)^2s^4-d(d-v^2)^2(1-e)^2.
\]
It is quadratic in the complicated detector expression $e$. The
representative $d$ is already produced by the circuit's square-factor
extraction; it requires no extra free variable. The new predicate still
has two retained rational parameters and one terminal algebraic root.
The degree-seven birational version in the companion twist note is
nonmonic in its last variable; converting it into the monic tower used
here restores the degree-nine product. Its smaller displayed degree is
therefore not by itself an additional determinant-degree improvement.

\section{A compatible four-contact scalar}
For odd $L\ge3$ put
\begin{align*}
 \beta&=1632X^2+2477,& A&=\beta(\beta-1),\\
 B_0&=(\beta-2401)(32\beta+343),&
 T_L&=\bigl(1224(67^{(L+1)/2}+1)\bigr)^2,\\
 q&=T_LA/B_0,&
 \rho_*&=102^2\,67^L\frac{7q}{(1+7q)^2}.
\end{align*}
Here $B_0$ denotes a contact polynomial, not the separate abc shift $B$.
The target remains $\tau=X^2+1$.

\begin{lemma}\label{scalarlemma}
For every rational $X$, $\rho_*$ is defined, positive, less than $1/167$,
and $v_{67}(\rho_*)=L$. Its negative valuations are even. For integral
$X$ it lies in $7\Q_2^{\times2}$, has $v_2(\rho_*)\ge10$,
$v_3(\rho_*)=6$, and $v_{17}(\rho_*)\ge4$. Moreover
\begin{equation}\label{height}
 h(\rho_*)\le4h(\beta)+2L\log67+66\log2
             <4h(\beta)+2L\log67+128.
\end{equation}
\end{lemma}
\begin{proof}
All four contact factors are positive because $\beta\ge2477$.
The old scalar $\rho=102^2\,67^Lq/(1+q)^2$ satisfies
$0<\rho<33/(144\cdot67)<1/167$, as $B_0/A<33$.
Also $q>T_L/33>1$, and
\[
 \frac{\rho_*}{\rho}=7\left(\frac{1+q}{1+7q}\right)^2<1.
\]
At integral 67-adic $X$, exact enumeration of the 67 residue classes
gives $AB_0(B_0+7\cdot1224^2A)\not\equiv0\pmod {67}$.
At negative $v_{67}(X)$ cancel the common valuation of $A,B_0$;
their ratio is $1/32$ modulo 67 and
$32+7\cdot1224^2\equiv22\pmod {67}$. Thus the valuation is exactly $L$.

For nonzero rational $u$, every negative valuation of $u/(1+u)^2$
is even: if $v_p(u)\ne0$ the resulting valuation is $|v_p(u)|$;
otherwise it is $-2v_p(1+u)$. Apply this to $u=7q$ and observe
that the remaining factors are a square and $67^L$, whose total
67-adic valuation has just been computed as positive. This argument
does not assert that $v_7(\rho_*)$ is positive; it may be negative even.

For integer $X$, $\beta\equiv13\pmod {32}$, so
$v_2(A)=v_2(B_0)=2$ and $AB_0/16\equiv3\pmod8$.
The old scalar was a 2-adic square; the ratio above changes its class
to seven. Since $v_2(T_L)\ge8$, its depth is at least ten.
The factor seven is a 3-adic square and the two normalizers are units
there, preserving depth six. At seventeen the depth remains at least
four, while the squareclass changes to seven. The old 17-adic square
requirement was used for the replaced level-17 pole converse; it is
not a hypothesis of BB Lemma~2.2 or of the deep formal-subgroup reference
choice, which uses continuity and valuation depth.

Finally $\|A\|_1=2$, $\|B_0\|_1=900064$, and
$T_L\le F_0\,67^L$, with $F_0=4\cdot1224^2\cdot67=401511168$.
After removing the factor $67^{2L}$, the numerator and denominator
coefficient one-norms of this degree-four rational map are bounded by
\[
 7\cdot102^2F_0\cdot2\cdot900064=52638002498271117312,
 \qquad(900064+14F_0)^2=31607518344686765056.
\]
Both are less than $2^{66}$. The elementary height inequality for a
rational map proves \eqref{height}.
\end{proof}

\section{Arithmetic transfer and retained thresholds}
In both general sources replace $\rho$ by $\rho_*$, and replace each
native-pole block by \eqref{product}, extracting $e=dk^2$ in the same
coefficient DAG. Leave the point charts, comparisons, contact locations,
weights, and unsquared scales as in v78. The finite exceptional set already
contains seven. Outside it, scalar squareclasses are multiplied by the
unit seven, so the parity of each consumed valuation is unchanged.
At a finite contact the old and new normalizers are units when $v_p(q)>0$;
when $v_p(q)<0$ the valuation of $q/(1+q)^2$ and that of
$7q/(1+7q)^2$ is $-v_p(q)$. Thus positive odd contact valuations and
their exact depths agree. The fixed 67-adic depth remains $L$.
The square-denominator implication of Lemma~\ref{pole} therefore supplies
every old local exclusion and divisibility statement.

The scalar contact radical, exceptional sets, positive scale properties
and target inequality $h(X^2+1)\le h(\beta)$ are unchanged.
The stronger estimate \eqref{height} implies the old $+128$ bound.
Consequently the entire effective three-scale ledger in
\cite[Sections 4--5]{scalar} applies verbatim, including the exact
inequalities
\[
 67^{287}\ge120\cdot2^{342},\qquad
 67^{287}-1\ge2^{1727},\qquad
 2^{1727-3\cdot551}12^6>18^3(2400+7\cdot1727)^6.
\]
They certify $L=143$ for the 27-variable source. For the two-scale source
the same old height contradiction has an ineffective constant; its
sufficient fixed odd $L$ remains ineffective.

For completeness first fix integral $X$ and the nonzero root label
$X^2+1$. Make the reference, Lucas-label, point and prime choices in
the source order. The actual pole inputs in this source are
\begin{equation}\label{actual}
 e_j=\frac{4\rho_*R_j}{\Tr(W_j^{-1})^2}.
\end{equation}
There is no extra $(1+\rho_*R_j)^2$ normalizer in \eqref{actual}.
At newly selected primes the successful-slot construction makes the trace
a unit and the input valuation positive odd. At the finitely many primes
in the scalar's support, the reference neighborhoods give the old
nonnegative odd valuations and preserve all even negative valuations.
Elsewhere the weight has even valuation and division by a square does
not change parity. Thus each $e_j$ has square reduced denominator.
At two, $R_j$ is close to one and the trace to two, so $e_j$ is in
$7\Q_2^{\times2}$ with depth at least ten. At the chosen real character
the existing neighborhoods give $R_j<2$ and $\Tr(W_j^{-1})>1$;
hence $0<e_j<8\rho_*<1$. These are precisely the hypotheses of
Lemma~\ref{pole}. Its density avoids the remaining nonidentically-zero
guards. This proves completeness with the same retained variables.

For the ordinary-abc source no scalar change is required. In the v78
notation its actual input is
\[
 \eta=7u\left(\frac{N(W)}{(1+u)W_0}\right)^2,
 \qquad d=7u.
\]
The existing completeness proof prepares $\eta$ in class seven with
square denominator, the required depth and $0<\eta<1$; the v80
published-theorem transfer explicitly checks the intrinsic-pole
preparation. Replace only its last radicand by
$(1-\eta)(v^4-d^2)(s^4-d^2)$. The twelve contact locations, first
height scale, even-valuation comparisons and balancing are unchanged.
The sufficient odd shift remains ineffective under ordinary abc.

\section{Balancing the squareclass representative}\label{sec:balanced}
The general-source degree bounds use a further rational coordinate
change. For any guarded nonzero factor $M$, replace an extraction
$e=d_{\rm old}k_{\rm old}^2$ by
\[
 d=d_{\rm old}M^2,\qquad k=k_{\rm old}/M.
\]
The product-pole charts are isomorphic by
$(v,s,Y)_{\rm new}=(Mv,Ms,M^4Y)_{\rm old}$. Thus existence, density,
terminal zero roots and the number of free variables are unchanged.
Every factor selected below was already required to be nonzero.

Here is a semantic specification of the selected normalization.
For detector $j$, write its scaled slope as the exact stored circuit
$a_j=N_j/D_j$, with $D_0=1$ and
$D_j=(B_0+7T_LA)^2$ for $j\ge1$. In order, take the ten contacts
\[
 (4j+2+\Theta,\ 4j+4+\Theta),\qquad
 \Theta=(0,-1/2,-1,1/2399,-32/407).
\]
For a contact $\theta_{j,\ell}=a/b$ in lowest terms, $b>0$, put
$L_{j,\ell}=bN_j-aD_j$. Exact independent circuit replay gives
\[
 d_{j,\rm old}^{-1}=C_jAB_0\prod_{\ell=1}^{10}L_{j,\ell},
\quad C_j=4879476\ (j\ge1),\quad
C_0=4879476(2\cdot2399\cdot407)^2.
\]
Choose
\[
 M_0=A\prod_{\ell=1}^{6}L_{0,\ell},\qquad
 M_j=A(\beta-2401)\prod_{\ell=1}^{5}L_{j,\ell}\quad(j\ge1).
\]
Consequently
\[
 d_0=\frac{A\prod_{\ell=1}^{6}L_{0,\ell}}
 {C_0B_0\prod_{\ell=7}^{10}L_{0,\ell}},\qquad
 d_j=\frac{A(\beta-2401)\prod_{\ell=1}^{5}L_{j,\ell}}
 {C_j(32\beta+343)\prod_{\ell=6}^{10}L_{j,\ell}}\quad(j\ge1).
\]
The numerator/denominator degree pairs are $(10,8)$ and $(71,67)$.
The total guard degrees remain 8215 and 14398. The complete coefficient
DAG fixes all constants and the choice of rational circuit, so the
formula does not depend on an arbitrary rescaling of $N_j,D_j$.
The finite search covered eighteen initial menus per route and all
fifty-five scalar/contact prefix menus; the selected combined choices
were then rebuilt and independently checked in both complete sources.
No global optimality of this factor redistribution is asserted.

\section{Zero strata, exact specifications and degree calculation}
The terminal-root policy is part of the new specification. One must
not silently impose nonzero terminal ordinates: at another algebraic
character the expression $e$ may equal one, making its radicand zero.
The final positive encoders of v78 allow all terminal zero roots.
Their inherited formal source-level ordinate guards are explicitly
removed, and the independent checker verifies that no denominator or
retained guard depends on such a root. The final polynomial is the
specified positive encoder, not the legacy guarded-source formula alone.

The point-prefix algebra, ordered point-guard recovery and positivity
guard deletions are unchanged. For one terminal root use its bounded
linear separator; for three use the audited cubic signed separator.
For two use the projective separator from Gabe Conant's notes with
Astra~6.2. The last construction gives numerator and denominator
operators $N_1,D_1$, with $D_1$ nonzero on every geometric character.
The equations are respectively the integral determinants
\[
 \det(I-YG\operatorname{Mult}(\Omega))=0,
 \qquad
 \det(\operatorname{Mult}(D_1)-YG\operatorname{Mult}(N_1))=0.
\]
The abc polynomial additionally has the source input factor $X$.
Multiplication on each Artin local factor has its residue character on
the diagonal, so nonreduced terminal specializations do not create
additional rational zeros. Rational-character recovery and the
subsequent arithmetic implications are those proved in the cited
source encoder papers, with the new terminal equations supplied above.

Here are the complete final degree data. For 25 variables the last
column is both the projective numerator and denominator degree;
in the other two rows it is the separator degree.
\begin{center}\begin{tabular}{rrrr}
\toprule
Variables&Algebra rank $n$&Guard degree $g$&Operator degree $w$\\\midrule
25&$2^{20}$&8215&37864\\
27&$2^{21}$&14398&73038\\
11&128&650&24830\\\bottomrule
\end{tabular}\end{center}
Consequently the bounds are $n(1+g+w)$ in the first two rows and
$1+n(1+g+w)$ in the third. The coefficient log-one-norm bounds
at the stored fixtures are respectively
\[
 6235064631296,\qquad28636742680576,\qquad12940032.
\]
For total degree $D$, input-factor indicator $\epsilon$ and $v$ total
variables, the support bound is the exact integer
\[
 {D-\epsilon+v\choose v}-{D-\epsilon-n+v-1\choose v},
\]
rounded upward to the displayed strict power of two.

The retained-coordinate counts are also explicit: the general source
has one input, eight root-chart parameters, ten auxiliary point
parameters, one slope, four or six pole parameters and one final
encoder variable. The abc source has one input, six point parameters,
one contact parameter, two pole parameters and one encoder variable.
Thus the counts are still 25, 27 and 11.

The package \path{computation/bhargava_burungale_v88} contains every
coefficient DAG, point/pole relation, denominator, guard, norm factor,
and final bound. The producer is rebuilt from the byte-pinned v78
archives. The independent checker does not import the new producer:
it reconstructs every changed scalar and pole, proves $dk^2=e$ exactly,
compares unchanged charts and tests with their sources, and replays
every operator majorant and guard deletion. A separate standard-library
checker verifies scalar residues, exceptional sets and effective
inequalities. These finite checks do not prove the imported infinite
arithmetic statements; their scope is recorded explicitly.

\section{The weak-$abc$ eleven-variable routes}
\label{sec:weak-bb}
Let $W_\kappa$ be the weak rational $abc$ inequality with exponent
$\kappa$, as in the v78 single-pole theorem \cite{weak}. Retain that theorem's
elliptic group and point-chart identification, one-point trace-height
bound, formal-comparison assertions, and quantified prime-pattern
selection with finite congruences, reduction orders and approximation.
The following modification replaces its native pole converse by the
stated class-seven case of Bhargava--Burungale Lemma~2.2 and Theorem~2.3.

\begin{theorem}
For each fixed $1<\kappa<3/2$, the v78 eleven-variable construction
admits the Bhargava--Burungale product pole without changing its contact
scalar, point charts, physical tests, or height inequality. In particular,
the following bounds hold for every sufficiently large fixed odd shift
$B$, under $W_\kappa$ and the retained premises just stated:
\[
\begin{array}{c|c|c|r|r|r}
\kappa&e&k&\text{old degree bound}&\text{new degree bound}&
\text{support bound}\\\hline
121/100&6&6&32542081&16460929&<2^{225}\\
16/15&4&5&19366273&9808001&<2^{218}
\end{array}
\]
Every count is eleven total rational coordinates, including the input.
The sufficient shift remains ineffective. The tabulated coefficient
bounds at $B=3$ are reproducibility fixtures, not certified defining
specializations.
\end{theorem}

\begin{proof}
Use the original notation
\[
 x=a-z,\quad y=a+z,\quad P=G(x),\quad Q=G(y),
 \quad J=(x-1)(y-1),\quad F=P^2+Q^3,
\]
\[
 u=\frac{J^{8(k-1)}PQF^2}
 {(P^3+Q^5)(P+Q^k)(P^k+Q)},\qquad
 W=\prod_{i=1}^6(1+(t_i/u)^2),\quad V=\operatorname{Tr}(W^{-1}),
\]
\[
 b=\frac{4u}{(1+u)^2V^2}.
\]
Keep every one of these rational expressions and all six physical tests.
Replace only the final pole input by $\epsilon=7b$. Choose a square
extraction $\epsilon=d_0t_0^2$ and the positive contact sum
$\delta=P+Q^k$. Put
\[
 d=d_0/\delta^2,\qquad t=t_0\delta,
 \qquad \epsilon=dt^2.
\]
The new two-parameter terminal quadratic equation is
\begin{equation}\label{eq:weak-bb-pole}
 Y^2=(1-\epsilon)(v^4-d^2)(s^4-d^2).
\end{equation}
The monic product-pole equation with input $\epsilon$ is recovered by
$v_{\mathrm{old}}=tv$, $s_{\mathrm{old}}=ts$, and
$Y_{\mathrm{old}}=t^4Y$; its radicand is $t^8$ times the right side of
\eqref{eq:weak-bb-pole}. Thus the scaling preserves rational solutions
and density on the existing nonzero input and square-multiplier locus.

For soundness, the old scalar-height argument supplies positive odd
valuation primes outside a fixed set $S$ containing $2,3,7,17$.
At such a prime, $v_p(\epsilon)=v_p(b)$ and
$v_p(d)\equiv v_p(b)\equiv v_p(u)\pmod2$. If $v_p(\epsilon)<0$ is odd,
then $v_p(1-\epsilon)$ is odd, whereas each factor $v^4-d^2$ or
$s^4-d^2$ has even valuation: its two summands have different valuations
modulo four. A zero factor cannot occur in this case, including when
$v=0$ or $s=0$. This contradicts \eqref{eq:weak-bb-pole}, even when
$Y=0$. Hence the same negative-odd-valuation obstruction holds.
All local saturation estimates of v78 therefore apply unchanged. In
particular, with $\nu=v_p(u)>0$ odd and the old $q_*\le0$, one obtains
$0<v_p(b)\le\nu+2q_*$ and then the original congruence $X\equiv n\pmod p$.
The scalar-height margins remain
\[
 \lambda=8374/14641>1/2,
 \qquad \lambda=99/160>1/2,
\]
respectively. The general condition for $1<\kappa<3/2$ is unchanged.

For completeness, retain the ordered integer witnesses of v78. Its
all-prime argument already proves that the reduced denominator of $b$
is a square, and at every fixed place its local preparation has
$b z^8\longrightarrow1$ as $z\longrightarrow\infty$.
Enlarge the temporary fixed preparation by the prime seven before
applying the prime-pattern theorem. Require $v_7(z)<0$ sufficiently
negative, so that $v_7(b)=-8v_7(z)$ is nonnegative and even. The product
$\epsilon=7b$ then has no negative odd valuation at seven; outside seven
this follows from the square denominator of $b$.
At two retain $b\in\mathbb Q_2^{\times2}$ and take
$v_2(b)\ge4$ even. Then $\epsilon$ has two-adic square class seven and
$1-\epsilon$ is a two-adic square. The real preparation can be strengthened
to $0<b<1/7$, giving $0<\epsilon<1$. These are finitely many prescribed
neighborhood conditions of precisely the type already included in the
retained selection and recurrence premise. They are imposed before
choosing the contact primes and do not change the order of choices.

Writing the positive squarefree parts of $\epsilon$ and $1-\epsilon$
as $\beta$ and $\eta$, respectively, gives
$\beta\equiv7\pmod8$, $\eta\equiv1\pmod8$, and
$\gcd(\beta,\eta)=1$. The local quadratic residue conditions follow
from $\epsilon+(1-\epsilon)=1$ and its square reduced denominator.
Bhargava--Burungale's stated class-seven theorem and the product-pole
parameterization therefore give rational witnesses whose $(v,s)$
projection is Zariski dense. The finite-character guard avoidance
argument transfers under the displayed nonzero rational scaling.

The old conic denominator $\tau^2+51d$ is absent. Accordingly, the final
compiler removes nine automatic guards, not the previous ten: the
root-ordinate and two auxiliary-abscissa factors before character
extraction, the root-abscissa factor afterward, four positive contact
sums, and the odd-weight normalizer $1+u$. The original balancing and
positive-product stages remain unchanged. The terminal root guard is
suppressed by the same final-root theorem; no earlier denominator or
retained guard depends on that root. Its zero stratum is included.

The unchanged three point blocks and the new terminal square block have
ranks $4,4,4,2$. The retained ten free coordinates and one encoder give
rank $128$ and eleven total variables. The independently reconstructed
pairs of guard and separator bounds are
\[
 (g,w)=(490,128110),\qquad (418,76206).
\]
Thus $\deg F_B\le1+128(1+g+w)$ gives the displayed degree bounds.
The same integer determinant has coefficient of $XY^0$ equal to one,
so its content is one. At $B=3$, the independently checked coefficient
one-norm bounds are
$2^{66654464}$ and $2^{42394624}$, respectively.
\end{proof}

The new checker reconstructs the unchanged contact and detector
expressions, every physical point/test comparison, the identity
$\epsilon=dt^2$, the product pole, all square extractions, all independent
operator and subalgebra-norm majorants, the terminal zero-root condition,
and the staged guard and determinant bounds. Its finite checks do not
prove the quantified arithmetic premises or an effective value of $B$.

\section{Variable-count limitations}
The product substitution lowers circuit degrees because it simplifies
dependence on a high-degree detector. It removes no free parameter.
The finite-support twist theorem in the companion note can share
parameters only when all actual inputs have one common squareclass.
That condition has not been proved for the present multiple detectors.
Nor does the published theorem give rational points on every proposed
one-parameter slice through an actual detector fiber.

The 163-, 179- and 56-variable quartic records, and the conditional
number-field counts, are therefore unchanged. The 21/9-variable targets
still require a global rational-realization theorem with every physical
lift and denominator condition. None of the new finite computations
establishes it.

\begin{thebibliography}{9}\raggedright
\bibitem{BB} M. Bhargava and A. Burungale,
\emph{A Diophantine definition of $\Z$ in $\Q$},
arXiv:2611.16130v2, 19 November 2026,
\url{https://arxiv.org/html/2611.16130v2}.
\bibitem{scalar} $\Z\Q$ Diophantine Project,
\emph{Four scalar contacts and unsquared comparison scales}, v78,
19 November 2026; complete source in this repository.
\bibitem{twelve} $\Z\Q$ Diophantine Project,
\emph{Twelve contacts at the first height scale}, v78,
19 November 2026; complete source in this repository.
\bibitem{weak} $\Z\Q$ Diophantine Project,
\emph{The single-pole eleven-coordinate construction under weak abc},
v78, 19 November 2026;
\path{papers/conditional_few_variables/weak11_audit_and_transfer_20261119_v78.tex}.
\bibitem{coauthor} Gabe Conant with Astra~6.2,
\emph{H10-quartic-constructions} and
\emph{H10-eleven-variable-constructions}, research bundles,
18 November 2026; archived with attribution and checksums in v78.
\end{thebibliography}
\end{document}
